\uuid{tesK}
\titre{Changements de variables/éléments simples}
\auteur{hanani ; mzari}
\organisation{exo7}
\licence{CC BY-NC-SA 4.0}
\datecreate{2020-01-20}
\niveau{L1}
\module{Analyse}
\chapitre{Calcul d'intégrales}
\sousChapitre{Autre}
\difficulte{3}
\theme{Calculs d'intégrales, changements de variables, éléments simples}
\exo7qcmid{lille-2.integrales.036}

\contenu{
\begin{question}
\qclassification{127.99}
\qtags{motcle=changements de variables/éléments simples}
Parmi les affirmations suivantes, cocher celles qui sont vraies :
\begin{answers}
\good{Le changement de variable $\displaystyle t=\sqrt{\frac{1+x}{1-x}}$ donne :
$$\displaystyle \int _0^{1/2}\sqrt{\frac{1+x}{1-x}}\mathrm{d}x=\int _1^{\sqrt{3}}\frac{4t^2\, \mathrm{d}t}{(t^2+1)^2}.$$}
\bad{$\displaystyle \int _1^{\sqrt{3}}\frac{\mathrm{d}t}{t^2+1}=\frac{\pi}{6}$.}
\good{$\displaystyle \int _1^{\sqrt{3}}\frac{\mathrm{d}t}{t^2+1}=\left[\frac{t}{1+t^2}\right]_1^{\sqrt{3}}+\int _1^{\sqrt{3}}\frac{2t^2\mathrm{d}t}{(t^2+1)^2}$.}
\bad{$\displaystyle \int _0^{1/2}\sqrt{\frac{1+x}{1-x}}\mathrm{d}x=\frac{\pi}{3}-\frac{\sqrt{3}}{2}+1$.}
\end{answers}
\vskip2mm
\begin{explanations}
Avec $\displaystyle t=\sqrt{\frac{1+x}{1-x}}$, on a : $\displaystyle x=\frac{t^2-1}{t^2+1}$. D'où $\displaystyle \mathrm{d}x=\frac{4t\, \mathrm{d}t}{(t^2+1)^2}$ et
$$\int _0^{1/2}\sqrt{\frac{1+x}{1-x}}\mathrm{d}x=\int _1^{\sqrt{3}}\frac{4t^2\, \mathrm{d}t}{(t^2+1)^2}.$$
Une intégration par parties avec $\displaystyle u=\frac{1}{t^2+1}$ et $v=t$ donne
$$\int _1^{\sqrt{3}}\frac{\mathrm{d}t}{t^2+1}=\left[\frac{t}{1+t^2}\right]_1^{\sqrt{3}}+\int _1^{\sqrt{3}}\frac{2t^2\mathrm{d}t}{(t^2+1)^2}.$$
Donc $\displaystyle \int _0^{1/2}\sqrt{\frac{1+x}{1-x}}\mathrm{d}x=2\int _1^{\sqrt{3}}\frac{\mathrm{d}t}{t^2+1}-2\left[\frac{t}{1+t^2}\right]_1^{\sqrt{3}}=\frac{\pi}{6}-\frac{\sqrt{3}}{2}+1$.
\end{explanations}
\end{question}
}
